\cotefile{83}
\scanpage{194}
% archive sheet number: 193

(fonctions \deleted{symétriques} \inserted{\unsure{déterminantes|déterminants|déterminées}} pour l'action du groupe $\unsure{\mathfrak S|\mathscr S}_{\{i,j,k\}}\subset\mathfrak S_4$). \unsure{Ça|Si} fait qu'en fait, il suffit comme générateurs les quatre quantités (en plus du quatuor $b_{,,},c_{,,},t_{..}$ du début)

\[
 t_{2,3,4},\,t_{1,3,4},\,t_{1,2,4},\,t_{1,2,3}
\]

Si $A$ est l'anneau engendré par les $b_{..},c_{..},t_{..}$, je présume que $B$ est fini libre de rang $2^4$ sur $A$, avec $8$ \unsure{génér.} \inserted{distincts $t_{ijk}$} sur $A$ soumis aux seules relations $(\ast)$ ci-dessus.

Mais il faudrait trouver la relation entre les $b_{\unsure{j|i}},c_j,t_{k,\ell}$ -- peut-être via les $t_{ijk}$. La clef pour la trouver \deleted{\unsure{est}} \inserted{sera} le calcul de l'expression

\[
 \operatorname{Tr} x_1x_2x_3x_4
\]

\deleted{\unsure{,}}
